\chapter*{Keywords and UNESCO Codes}
\addcontentsline{toc}{chapter}{Keywords and UNESCO Codes}
\markboth{Keywords and UNESCO Codes}{Keywords and UNESCO Codes}

\vspace{0.4cm}
\noindent\textbf{Keywords}
\vspace{0.3cm}

\noindent
Discrete-event simulation \quad$\cdot$\quad Warehouse operations \quad$\cdot$\quad
Workforce capacity planning

\vspace{0.15cm}
\noindent
Order sequencing \quad$\cdot$\quad Service-level agreements \quad$\cdot$\quad
Decision support

\vspace{1.2cm}
\noindent\textbf{UNESCO Codes}
\vspace{0.3cm}

\noindent
\begin{tabularx}{\textwidth}{llX}
\toprule
\textbf{Code} & \textbf{Field (original)} & \textbf{Field (English)} \\
\midrule
1203.26 & Simulaci\'on & Simulation \\
1207.12 & Colas & Queueing \\
5311.07 & Investigaci\'on operativa & Operations research \\
5311.09 & Organizaci\'on de la producci\'on & Production organisation \\
\bottomrule
\end{tabularx}

\vspace{0.6cm}
\noindent\footnotesize
Codes follow the UNESCO International Nomenclature for Fields of Science and Technology
(\emph{Nomenclatura Internacional de la UNESCO para los campos de la Ciencia y la
Tecnolog\'ia}), six-digit level.
